% ============================================================
%  Esercizi — Appendice
%  Livello: prima media (prima parte dell'anno)
%  Argomenti: quattro operazioni, potenze, geometria piana
% ============================================================

\chapter{Esercizi}
\section{Esercizi}

% --- Espressioni aritmetiche -----------------------------------
\subsection{Espressioni aritmetiche}

Calcola il valore delle seguenti espressioni, rispettando le regole
di precedenza.

\begin{enumerate}[label=\textbf{\arabic*.}, leftmargin=2em, itemsep=1.2ex]

      % --- Solo quattro operazioni ---

      \item $48 : 6 + 5 \cdot 3 - 7$

      \item $(120 - 45 : 9) \cdot 2 - 13$

      \item $\{[72 : (3 + 6) - 4] \cdot 5 + 8\} : 9$

            % --- Con le potenze ---

      \item $3^{4} - (2^{3} + 10) \cdot 2$

      \item $\{5^{2} \cdot [3 \cdot (4^{2} - 7) + 1] - 10\} : 5$

\end{enumerate}

% --- Problemi --------------------------------------------------
\subsection*{Problemi}

\begin{enumerate}[label=\textbf{\arabic*.}, leftmargin=2em, itemsep=1.2ex,
            start=6]

      % --- Problema di aritmetica ---

      \item \textbf{La gita scolastica.}

            Una classe di $24$ allievi organizza una gita. Il costo del pullman
            è di $312$ franchi e il biglietto di ingresso al museo è di
            $7$ franchi per allievo. Ogni allievo versa $22$ franchi.

            \begin{enumerate}[label=\alph*), leftmargin=1.5em, itemsep=0.6ex]
                  \item Qual è il costo totale della gita?
                  \item Quanti franchi avanzano alla classe dopo aver pagato tutte
                        le spese?
            \end{enumerate}

            % --- Problema di geometria ---

      \item \textbf{Il giardino rettangolare.}

            Un giardino ha forma rettangolare. Il lato più lungo misura $15$ m
            e il lato più corto misura $9$ m.

            \begin{enumerate}[label=\alph*), leftmargin=1.5em, itemsep=0.6ex]
                  \item Calcola il perimetro del giardino.
                  \item Calcola l'area del giardino.
                  \item Il proprietario vuole recintare il giardino con una staccionata.
                        Il materiale costa $8$ franchi al metro lineare.
                        Qual è il costo totale della recinzione?
            \end{enumerate}

\end{enumerate}

% --- Soluzioni -------------------------------------------------
\subsection*{Soluzioni}

\begin{enumerate}[label=\textbf{\arabic*.}, leftmargin=2em, itemsep=0.8ex]

      \item $48 : 6 + 5 \cdot 3 - 7 = 8 + 15 - 7 = \mathbf{16}$

      \item $(120 - 45 : 9) \cdot 2 - 13
                  = (120 - 5) \cdot 2 - 13
                  = 115 \cdot 2 - 13
                  = 230 - 13
                  = \mathbf{217}$

      \item $\{[72 : (3 + 6) - 4] \cdot 5 + 8\} : 9
                  = \{[72 : 9 - 4] \cdot 5 + 8\} : 9
                  = \{[8 - 4] \cdot 5 + 8\} : 9
                  = \{4 \cdot 5 + 8\} : 9
                  = \{20 + 8\} : 9
                  = 28 : 9$

            \quad\textit{(corretto: non è divisione esatta;
                  si lascia in forma $28 : 9$ oppure $\approx 3{,}1$
                  se introdotta l'approssimazione)}

      \item $3^{4} - (2^{3} + 10) \cdot 2
                  = 81 - (8 + 10) \cdot 2
                  = 81 - 18 \cdot 2
                  = 81 - 36
                  = \mathbf{45}$

      \item $\{5^{2} \cdot [3 \cdot (4^{2} - 7) + 1] - 10\} : 5$

            \begin{align*}
                   & = \{25 \cdot [3 \cdot (16 - 7) + 1] - 10\} : 5 \\
                   & = \{25 \cdot [3 \cdot 9 + 1] - 10\} : 5        \\
                   & = \{25 \cdot [27 + 1] - 10\} : 5               \\
                   & = \{25 \cdot 28 - 10\} : 5                     \\
                   & = \{700 - 10\} : 5                             \\
                   & = 690 : 5                                      \\
                   & = \mathbf{138}
            \end{align*}

      \item \textbf{La gita scolastica.}

            \begin{enumerate}[label=\alph*), leftmargin=1.5em]
                  \item Costo totale:
                        $312 + 24 \cdot 7 = 312 + 168 = \mathbf{480}$ fr.
                  \item Versamento totale: $24 \cdot 22 = 528$ fr.
                        Resto: $528 - 480 = \mathbf{48}$ fr.
            \end{enumerate}

      \item \textbf{Il giardino rettangolare.}

            \begin{enumerate}[label=\alph*), leftmargin=1.5em]
                  \item Perimetro: $2 \cdot (15 + 9) = 2 \cdot 24 = \mathbf{48}$ m
                  \item Area: $15 \cdot 9 = \mathbf{135}$ m\textsuperscript{2}
                  \item Costo recinzione: $48 \cdot 8 = \mathbf{384}$ fr.
            \end{enumerate}

\end{enumerate}